% =====================================================================
\section{Kết luận}
\label{sec:conclusion}
% =====================================================================

\subsection{Kết luận chính}
\label{sec:concl-main}

Báo cáo đã xây dựng và chạy một quy trình thực nghiệm đóng kín cho đề tài M8: một phép chia
benchmark được đóng băng với các tập nhãn lồng nhau, một mô hình ResNet-18 hai đầu ra cho ảnh
$32 \times 32$, ba cấu hình đối chứng có thể tách bạch phần nào hai yếu tố \enquote{thêm nhiệm
vụ phụ} và \enquote{mở rộng nguồn ảnh cho nhiệm vụ phụ}, cùng các phép đo chi phí huấn luyện và
suy luận.

Khối thực nghiệm chính hoàn thành \RRuns{} lượt chạy trong \RTotalHours{} giờ. Kết quả trên tập
test cho thấy \RFindingSentence

Về mặt chi phí, mô hình suy luận của hai cấu hình \cfg{E1} và \cfg{E2} là \emph{giống nhau} với
\RParamInference{} tham số (\RParamInferenceMiB{} MiB ở FP32); Rotation Head chỉ thêm
\RParamRotationHead{} tham số trong giai đoạn huấn luyện. Do đó mọi khác biệt về chất lượng
phân loại là do chất lượng đặc trưng học được, không do thay đổi kiến trúc suy luận.

\subsection{Kết luận về phần mở rộng contrastive}
\label{sec:concl-ext}

Khối mở rộng \cfg{E3}--\cfg{E6} đã được triển khai đầy đủ theo bảng 5 của đề cương (hai hàm
mất mát NT-Xent và SupCon, Projection Head chỉ dùng khi huấn luyện) và chạy với siêu tham số
$\lambda_c$ khóa bằng validation. \RExtensionConclusion

\subsection{Tiêu chí hoàn thành đã đối chiếu}
\label{sec:concl-criteria}

\begin{table}[htbp]
  \centering
  \small
  \begin{tabular}{@{}p{7.4cm} p{7.6cm}@{}}
    \toprule
    Tiêu chí & Trạng thái đối chiếu \\
    \midrule
    Xây dựng đúng quy trình đã khai báo &
    Đạt. Ba cấu hình, ba mức nhãn, ba lần lặp, 100 epoch, $\lamrot = 0{,}5$; xem
    \tabref{tab:hyperparams} và \tabref{tab:per-run}. \\
    \addlinespace
    Kiểm tra được việc tách dữ liệu và che nhãn &
    Đạt. \code{results/benchmark\_audit.json} xác nhận không trùng split, các tập $\DL$ lồng
    nhau, $\DL \cap \DU = \varnothing$ và $\DL \cup \DU$ không giao với $\Dval$ và $\Dtest$. \\
    \addlinespace
    Thực hiện đối chứng bắt buộc &
    Đạt. \cfg{E1}, \cfg{E2-L} và \cfg{E2} đều được chạy đủ ở ba mức nhãn và ba lần lặp. \\
    \addlinespace
    Báo cáo trung thực các lượt chạy &
    Đạt. Mọi lượt đều được đưa vào báo cáo; không loại bỏ lượt nào; kết quả âm được phân tích. \\
    \addlinespace
    Cung cấp mã nguồn tái lập &
    Đạt. Xem \secref{sec:appendix-repro}. \\
    \addlinespace
    Mở rộng contrastive (tùy chọn) &
    Đạt. \cfg{E3}--\cfg{E6} triển khai đủ và có kết quả; xem \secref{sec:results-ext}. \\
    \bottomrule
  \end{tabular}
  \caption{Đối chiếu tiêu chí hoàn thành của đề cương với trạng thái thực tế của báo cáo.}
  \label{tab:completion}
\end{table}

\subsection{Giới hạn của kết luận}
\label{sec:concl-limits}

Kết luận của báo cáo chỉ áp dụng cho \emph{phép gộp hai siêu lớp} do đề tài quy ước, các tỷ lệ
nhãn và giao thức đã khảo sát, và cho kiến trúc cụ thể đã dùng. Cụ thể:

\begin{itemize}[leftmargin=1.2cm,itemsep=3pt]
  \item \textbf{Không} so sánh trực tiếp Accuracy hai lớp với bảng xếp hạng CIFAR-10 mười lớp.
  \item \textbf{Không} khẳng định đặc trưng học được tổng quát tốt hơn các phương pháp học tự
        giám sát khác, vì không có thực nghiệm đối chứng với chúng.
  \item \textbf{Không} khẳng định thuật toán là mới: đóng góp là một quy trình đánh giá có thể
        tái lập, không phải một kiến trúc.
  \item \textbf{Không} khẳng định tối ưu năng lượng: không có phép đo năng lượng nào được thực
        hiện.
  \item \textbf{Không} khẳng định đã xây dựng hệ thống IoT hoàn chỉnh: không có thử nghiệm nào
        trên thiết bị IoT thật.
  \item Ba lần lặp chỉ cho bằng chứng ban đầu về độ ổn định; không có kết luận nào về ý nghĩa
        thống kê được đưa ra.
\end{itemize}

\subsection{Hướng phát triển}
\label{sec:concl-future}

Ngắn hạn, ưu tiên tăng số lần lặp ở mức $10\,\%$ nhãn và triển khai khảo sát $\lamrot$ đã khai
báo. Trung hạn, việc tách thống kê BatchNorm theo nhánh sẽ cho phép quy kết hiệu ứng rõ ràng
hơn. Dài hạn, nếu mục tiêu là thiết bị biên, cần thay backbone bằng một kiến trúc hiệu quả tham
số và thực hiện lượng tử hóa có hiệu chỉnh, sau đó đo lại độ trễ và độ chính xác trên thiết bị
đích thay vì trên máy tính.
